\chapter{Annexe : Structure physique d'une ruche}
%==============================================================================
Pour mémoire, une ruche à cadres amovibles se compose, de bas en haut, des
éléments suivants : plancher / planche d'envol, corps (socle) contenant les
cadres de couvain, grille à reine, une ou plusieurs hausses contenant les cadres
de récolte du miel, couvre-cadre et toit. Cette structure justifie les règles de
composition (corps obligatoire, jusqu'à 5 hausses, 1 à 10 cadres par élément)
retenues dans le modèle métier.

\begin{figure}[htbp]
\centering
\includegraphics[width=0.55\linewidth]{hive-000}
\caption{Vue éclatée d'une ruche à cadres : toit, hausses, corps, plancher et planche d'envol.}
\label{fig:eclate-ruche}
\end{figure}

